% compendium/ml-alternative-data/labelling-sample-weights.tex
% Entries migrated on 2026-10-16 from the legacy AFML labelling and sample-weights files (slice 32,
% legacy/ml-alternative-data/data/datalabel.tex and datasampleweights.tex): the methods owned by
% statistical-learning-foundations, the time-validity pitfall by financial-ml-evaluation. Dispositions in
% notes/migration-slices.md; the files verbatim in notes/artefacts/statistical-learning-foundations/.

\section{Labels and sample weights for overlapping financial outcomes}

\begin{qfmethod}[Triple-barrier labelling: the fixed-horizon method and its limits, volatility-scaled barriers, first touch, barrier configurations, and labelling side against size]\label{meth:triple-barrier-labelling}
\qffield{Statement} \emph{Fixed-horizon labelling} gives an observation $X_i$, drawn from bars $t=1,\dots,T$, the label $y_i=-1$, $0$ or $1$ as the return $r_{t_{i,0},t_{i,0}+h}=p_{t_{i,0}+h}/p_{t_{i,0}}-1$ is below $-\tau$, within $\pm\tau$, or above $\tau$, with $t_{i,0}$ ``the index of the bar immediately after $X_i$ takes place'' and $\tau$ a constant \citep[chap.~3, sec.~3.2, p.~43]{alpha-afml}. López de Prado gives three reasons to avoid it: time bars have poor statistical properties; one $\tau$ is applied whatever the volatility, so a $\tau$ of 1 percent labels most observations 0 at night and few at the open; and it ignores the path, labelling as winners positions that ``would have been stopped-out'' (p.~44). The \emph{triple-barrier method} labels an observation by the first of three barriers touched: an upper (profit-taking) and a lower (stop-loss) horizontal barrier, set as multiples of a volatility estimate, and a vertical barrier, an expiration limit counted from the position's start; the upper gives $1$, the lower $-1$, and the vertical either the sign of the return or $0$, López de Prado preferring the sign (sec.~3.4, p.~45). With $t_{i,1}\le t_{i,0}+h$ the first-touch time, the label's return is $r_{t_{i,0},t_{i,1}}$, which uses the prices at $t_{i,0}$ and $t_{i,1}$ and compounds the one-bar returns $r_{t-1,t}$ with $t_{i,0}<t\le t_{i,1}$; the barriers need not be symmetric (p.~45).

\qffield{Conditions} \emph{Volatility.} The barrier width is a target such as an exponentially weighted standard deviation of daily returns (Snippet~3.1, default span 100 days, pp.~44--45), computed from data up to the event, so it is known at the decision time; observations whose target is below a minimum return are skipped (Snippet~3.3, p.~48). \emph{Configurations.} Writing $[pt,sl,t1]$ for which barriers are active, López de Prado calls $[1,1,1]$, $[0,1,1]$ and $[1,1,0]$ useful (the last ``somewhat unrealistic'', holding as long as it takes); $[0,0,1]$, $[1,0,1]$ and $[1,0,0]$ less realistic ($[0,0,1]$ is the fixed-horizon method, still useful on activity bars; $[1,0,0]$ ``could mean being locked on a losing position for years''); and $[0,1,0]$ and $[0,0,0]$ illogical, the last producing no label (p.~46). \emph{Side and size.} Without a model for the side of the bet, the profit-taking and stop-loss barriers cannot be told apart, so ``either there are no horizontal barriers or \dots\ the horizontal barriers must be symmetric'' (sec.~3.5, p.~48); with a known side, the barriers may be asymmetric and the label becomes a meta-label (\hyperref[meth:meta-labelling]{meta-labelling}).

\qffield{Practice} López de Prado's code compares bar closes with the barriers using strict inequalities and takes the earliest of the three touch times (Snippets~3.2--3.3, pp.~45--48); the label's return is measured to the close at $t_{i,1}$, so it can overshoot the barrier. His vertical barrier is set in calendar days (Snippet~3.4, p.~49). Interpretation, checked on toy paths (\nolinkurl{afml-checks.py}): with closes alone a bar cannot touch both horizontal barriers, but with highs and lows one bar can, and the text gives no rule for which came first; an implementation must choose one (stop-loss first, finer data, or drop the label) and record it. The code also starts the path at the event's own timestamp, whereas the definition starts it at the next bar: the entry price should be the first one tradable after the features are known. \emph{Arbitrage labels} (slice~36, from the legacy feature-engineering file; a routing exception from \texttt{lob-deep-learning}). The legacy principle that a spread which later converges is not automatically a good label, because the interim drawdown may breach margin, borrow, inventory or risk limits, so that labels should record both terminal convergence and maximum adverse excursion, is this method's path argument applied to a spread: a fixed-horizon label ignores positions that ``would have been stopped-out'' (p.~44), and a stop-loss barrier set at the drawdown the capital can bear makes the label path-dependent (Interpretation; Harris's staying power in \hyperref[res:arbitrage-risks]{arbitrage risks}). The legacy event labels (entry when the spread crosses an estimated arbitrage bound after costs; confirmation from flow; failure on specific news, borrow loss, deal break or hedge-ratio change; exit on convergence, stop-out, forced liquidation, expiry or capital reallocation) map onto the events that seed the barriers (\hyperref[meth:event-based-sampling]{event-based sampling}) and onto the three barriers; ``confirmation'' and ``failure'' are the legacy notes' categories, not López de Prado's, and the arbitrage bound is \hyperref[def:arbitrage]{arbitrage}'s. The statistical-arbitrage pathway, which would test such labels, is not open.

\qffield{Pitfalls} The legacy notes called the vertical barrier ``physical'', the $[1,0,0]$ position ``loose'' and the $[1,0,1]$ losses ``immediate'' rather than intermediate, dropped the path limitation of fixed horizons, the preference for the sign at the vertical barrier and the side-versus-size rule. A label uses the prices from $t_{i,0}$ to $t_{i,1}$: labels near the end of a training sample use prices from the test period; see \hyperref[pit:label-weight-leakage]{time-validity of labels and weights}.

\qffield{See also} \hyperref[meth:event-based-sampling]{event-based sampling} (the events that seed the barriers); \hyperref[meth:meta-labelling]{meta-labelling}; \hyperref[def:label-concurrency-uniqueness]{label concurrency and uniqueness}; \hyperref[meth:sample-weights-attribution-decay]{sample weights}.

\qffield{Further reading} López de Prado, chap.~10 (bet sizing from predicted probabilities), not read. The legacy file is kept in \nolinkurl{notes/artefacts/statistical-learning-foundations/}. The arbitrage-label blocks merged in slice~36 are kept in \nolinkurl{notes/artefacts/lob-deep-learning/}.

\qffield{Index} labelling; fixed-time horizon method; triple-barrier method; profit taking; stop loss; vertical barrier; expiration limit; first touch; dynamic thresholds; daily volatility estimate; barrier configurations; side and size; arbitrage labels; maximum adverse excursion

\qffield{Status} Verification: partly checked 2026-10-16 (against the text layer and page images of a local copy of \texttt{alpha-afml}, chap.~3, pp.~43--50, Snippets~3.1--3.5; toy paths labelled; the returns a label uses restated 2026-10-16; slice~36: the merged arbitrage-label blocks compared with p.~44). Stability: draft. Currency: timeless.
\end{qfmethod}

\begin{qfmethod}[Meta-labelling: a secondary model that decides whether to act on a primary model's side and how much to bet, with its claimed benefits stated as claims]\label{meth:meta-labelling}
\qffield{Statement} When a primary model (any model: ML, econometric, technical, fundamental or a human) sets the side of a bet, \emph{meta-labelling} trains a secondary model only to decide its size, including zero \citep[chap.~3, secs.~3.6 and~3.8, pp.~50 and~53--54]{alpha-afml}. Triple-barrier events are labelled with the primary side $s_i\in\{-1,1\}$: the meta-label is $1$ if the side-adjusted return $s_i\,r_{t_{i,0},t_{i,1}}$ is positive and $0$ otherwise, and the horizontal barriers may be asymmetric because profit-taking and stop-loss are now distinguishable (Snippets~3.6--3.7, pp.~50--51). The secondary model predicts the meta-label; when it predicts $1$, ``we can use the probability of this secondary prediction to derive the size of the bet, where the side (sign) of the position has been set by the primary model'' (p.~51).

\qffield{Conditions} López de Prado's recipe is to build a primary model with high recall, even at low precision, and to apply meta-labelling to its predicted positives so as to filter out false positives (sec.~3.7, pp.~52--53). He states that meta-labelling ``is particularly helpful when you want to achieve higher F1-scores'' and ``will increase your F1-score'' (pp.~52--53), that it lets ML sit on top of a white-box model, that ``the effects of overfitting are limited'' because ML decides only the size, that it allows separate long and short models, and that it concentrates ML on sizing (p.~53). These are claims without evidence in the chapter; his Exercise~3.5 asks the reader to compare the primary and secondary models' precision, recall and F1 (p.~55). Interpretation, derived here: filtering predicted positives can only keep or reduce recall, so F1 rises only if the precision gained outweighs the recall lost; it is not guaranteed.

\qffield{Practice} Interpretation: the meta-labels must be built from the primary model's \emph{out-of-sample} predictions (for example walk-forward, or cross-fitted with purging), because in-sample primary signals carry the primary model's overfitted hit rate, which the secondary model would learn. López de Prado's quantamental recipe says instead to ``fit your model again on the same training set, but this time using the meta-labels'' (sec.~3.8, p.~53), with no out-of-sample condition. For a discretionary portfolio manager he suggests features ranging from market information to stress and sleep (p.~54). Under heavy class imbalance he drops rare labels recursively (Snippet~3.8, p.~54) or uses balanced class weights (\hyperref[meth:sample-weights-attribution-decay]{sample weights}).

\qffield{Pitfalls} The legacy notes' steps (train a binary primary model, choose a high-recall threshold, possibly from ROC curves, concatenate its features and predictions, and call the final prediction positive ``only where both are true''), and their pipeline figure, follow a Hudson \& Thames blog example, not López de Prado \citep[``Model Architecture'' section]{ht-meta-labeling-toy}. In López de Prado the primary side is $\pm1$ and the secondary output sets a size; ``both true'' is the special case of sizing with a 0/1 threshold. The legacy's five feature groups for the secondary model, beyond the concatenation, and its three limitations of meta-labelling were found in neither source; the first (``if model has overfit the data, meta-labelling will not add much value'') contradicts López de Prado's own claim, and the third (it trades recall for precision and needs many trades) is consistent with his precision--recall trade-off (p.~52). They are not adopted, and kept in the artefact.

\qffield{See also} \hyperref[meth:triple-barrier-labelling]{triple-barrier labelling}; \hyperref[pit:label-weight-leakage]{time-validity of labels and weights}; \hyperref[pit:parameter-overfitting]{parameter overfitting}; \hyperref[meth:alpha-signal-combination]{alpha signal combination}.

\qffield{Further reading} Evidence on whether meta-labelling improves after-cost results is a backlog item of the pathway. The legacy file and its TikZ figure are kept in the pathway's artefact folder, \nolinkurl{notes/artefacts/} \nolinkurl{statistical-learning-foundations/}.

\qffield{Index} meta-labelling; primary model; secondary model; side and size; bet sizing; precision; recall; F1 score; false positives; confusion matrix; quantamental; class imbalance

\qffield{Status} Verification: partly checked 2026-10-16 (against the text layer and page images of a local copy of \texttt{alpha-afml}, chap.~3, pp.~50--55; the Hudson \& Thames page read for provenance only). Stability: draft. Currency: timeless; the Hudson \& Thames page as retrieved 2026-10-16.
\end{qfmethod}

\begin{qfdefinition}[Overlapping outcomes, label concurrency and average uniqueness]\label{def:label-concurrency-uniqueness}
\qffield{Statement} A label $y_i$ is a function of the prices over $[t_{i,0},t_{i,1}]$; when $t_{i,1}>t_{j,0}$ for $i<j$, $y_i$ and $y_j$ ``will both depend on a common return'', and the labels are not IID whenever some consecutive outcomes overlap, $t_{i,1}>t_{i+1,0}$ \citep[chap.~4, sec.~4.2, p.~59]{alpha-afml}. Restricting horizons to remove overlaps ties the sampling frequency to the label horizon, or, with path-dependent labels, to the first barrier touched, which López de Prado calls ``a terrible solution'' (pp.~59--60). Two labels are \emph{concurrent} at $t$ when both depend on the return $r_{t-1,t}=p_t/p_{t-1}-1$ (sec.~4.3, p.~60). With returns indexed by their ending time $t$, a label from $t_{i,0}$ to $t_{i,1}$ depends on exactly the returns with $t_{i,0}<t\le t_{i,1}$, so
\[ 1_{t,i}=\begin{cases}1&t_{i,0}<t\le t_{i,1},\\0&\text{otherwise,}\end{cases} \]
the \emph{concurrency} is $c_t=\sum_i1_{t,i}$, the \emph{uniqueness} of label $i$ at $t$ is $u_{t,i}=1_{t,i}/c_t$, and its \emph{average uniqueness} is
\[ \bar u_i=\frac{\sum_tu_{t,i}}{\sum_t1_{t,i}}, \]
the reciprocal of the harmonic mean of $c_t$ over the label's life, between 0 and 1 (sec.~4.4, p.~61).

\qffield{Conditions} $\bar u_i$ ``requires information that is not available until a future time'', the end time $t_{i,1}$; López de Prado says this is not a problem because the $\bar u_i$ ``are used on the training set in combination with label information, and not on the testing set'', and are not used to forecast the label (p.~62). The no-leakage claim therefore holds only under that use; see \hyperref[pit:label-weight-leakage]{time-validity of labels and weights}. \emph{Index conventions.} López de Prado writes $1_{t,i}=1$ ``if and only if $[t_{i,0},t_{i,1}]$ overlaps with $[t-1,t]$'' (p.~60); read with closed intervals that may touch at an endpoint, that would also mark the return ending at $t_{i,0}$ and the one starting at $t_{i,1}$, which his own toy matrix (p.~65) does not, so the rule above is the positive-length overlap. Any indexing that marks the same returns is equivalent: indexing returns by their starting time $s$, as $r_{s,s+1}$, a label marks $t_{i,0}\le s<t_{i,1}$. López de Prado's toy example (labels on $r_{0,3}$, $r_{2,4}$ and $r_{4,6}$, sec.~4.5.3, p.~65) is consistent in this sense: its printed matrix has rows for the returns ending at $t=1,\dots,6$, and Snippet~4.6 codes it as spans $\{0{:}2,\,2{:}3,\,4{:}5\}$, the start-time indexing with each end shifted down by one; those coded values are not the labels' end times. His production code is a different convention: Snippets~4.1 and~4.3 mark every bar from the event's timestamp $t_{i,0}$ to $t_{i,1}$ inclusive (\texttt{count.loc[tIn:tOut]}, \texttt{indM.loc[t0:t1]}), and Snippet~4.10 sums \texttt{ret.loc[tIn:tOut]} with \texttt{ret} the log return ending at each bar (pp.~60--61, 64 and~69). Applied to event times and returns indexed by ending time, that also counts the return ending at $t_{i,0}$, from before the label starts, so a label that starts where another ends acquires a return it never uses and the two appear concurrent. Including the start timestamp is therefore not an equivalent reindexing; the inclusive code is right only when its inputs are already shifted as in Snippet~4.6.

\qffield{Practice} Checked on the toy example (\nolinkurl{afml-checks.py}, and by brute force over the returns each label compounds in \nolinkurl{slice32-followup-checks.py}): with the rule above, concurrency $(1,1,2,1,1,1)$ for $t=1,\dots,6$ and average uniqueness $(0.833,0.75,1)$, López de Prado's matrix and values (pp.~65--66). The inclusive rule applied to the same events, rows $t=0,\dots,6$, gives concurrency $(1,1,2,2,2,1,1)$ and average uniqueness $(0.75,0.5,0.833)$: labels~2 and~3, adjacent at $t=4$, then share the return $r_{3,4}$, which label~3 does not use. Interpretation: the mean of $\bar u_i$ over the sample is the number of effectively independent outcomes per observation, the quantity behind the bootstrap and bagging choices in \hyperref[meth:sequential-bootstrap]{the sequential bootstrap}.

\qffield{Pitfalls} The legacy notes wrote $r_{t-1,t}=p_t/p_{t-1}=1$ for $-1$, and dropped the training-set condition from ``hence there is no information leakage''.

\qffield{See also} \hyperref[meth:triple-barrier-labelling]{triple-barrier labelling}; \hyperref[meth:sequential-bootstrap]{sequential bootstrap}; \hyperref[meth:sample-weights-attribution-decay]{sample weights}.

\qffield{Further reading} The legacy file is kept in \nolinkurl{notes/artefacts/statistical-learning-foundations/}.

\qffield{Index} overlapping outcomes; non-IID labels; concurrent labels; concurrency; uniqueness; average uniqueness; label span; indicator matrix

\qffield{Status} Verification: partly checked 2026-10-16 (against the text layer and page images of a local copy of \texttt{alpha-afml}, chap.~4, pp.~59--62 and 64--66; the toy example recomputed; the indicator convention restated and the inclusive code convention computed 2026-10-16, pp.~60--61, 64--65 and~69). Stability: draft. Currency: timeless.
\end{qfdefinition}

\begin{qfmethod}[Bootstrapping overlapping labels: redundancy, oversampling, and the sequential bootstrap]\label{meth:sequential-bootstrap}
\qffield{Statement} Drawing $I$ times with replacement from $I$ items misses a given item with probability $(1-I^{-1})^I\to e^{-1}$, so about $1-e^{-1}\approx0.63$ of the items are drawn \citep[chap.~4, sec.~4.5, p.~62]{alpha-afml}. When the mean average uniqueness $I^{-1}\sum_i\bar u_i\ll1$, bootstrapped in-bag observations are redundant with one another and similar to out-of-bag ones: bagged trees become ``very similar copies of a single overfit decision tree'' and out-of-bag accuracy ``will be grossly inflated'' (pp.~62--63). Dropping overlapping outcomes loses too much information; better is to sample only a fraction near the mean uniqueness (for example \texttt{max\_samples} in scikit-learn's bagging classifier) (p.~63). The \emph{sequential bootstrap} draws with changing probabilities: the first draw is uniform, $\delta_i^{(1)}=I^{-1}$; after the draws $\varphi$ so far, label $j$'s uniqueness given those draws is $u^{(2)}_{t,j}=1_{t,j}(1+\sum_{k\in\varphi}1_{t,k})^{-1}$, its average $\bar u^{(2)}_j=\sum_tu^{(2)}_{t,j}/\sum_t1_{t,j}$, and the next draw has probabilities
\[ \delta^{(2)}_j=\frac{\bar u^{(2)}_j}{\sum_{k=1}^I\bar u^{(2)}_k}, \]
updated after every draw until $I$ draws are made; repeats remain possible but less likely (sec.~4.5.1, pp.~63--64).

\qffield{Conditions} López de Prado says the sequential sample ``will be much closer to IID'', which ``can be verified by measuring an increase'' in mean uniqueness (p.~64). His evidence is a Monte Carlo on random spans (10 observations, 100 bars, label spans of 1 to 4 bars from Snippet~4.7's \texttt{maxH=5}, $10^6$ runs): median mean uniqueness 0.6 for the standard bootstrap against 0.7 for the sequential one, with an ANOVA test of the difference in means (Figure~4.2, pp.~66--68). That shows higher average uniqueness on those synthetic spans, a proxy for less redundancy; it is not a guarantee for every sample, nor evidence of better out-of-sample accuracy. The cost is one pass over all labels per draw.

\qffield{Practice} Checked (\nolinkurl{afml-checks.py}): on López de Prado's toy example, after label 2 is drawn the probabilities are $(5,3,6)/14$, as printed (p.~66), and after labels 2 and 3 they are $(5,3,3)/11$ (his Exercise~4.6). A 400-run Monte Carlo with his settings gives medians 0.600 and 0.700, with the sequential sample less unique than the standard one in 22 percent of runs (an independent 3,000-run check found 23 percent).

\qffield{Pitfalls} López de Prado writes that with at most $K\le I$ non-overlapping outcomes the probability of missing an item is $(1-K^{-1})^I$, ``approximated as $(1-I^{-1})^{I\frac{K}{I}}\approx e^{-K/I}$'', so that the share of unique observations drawn is $1-e^{-K/I}\le1-e^{-1}$ (p.~62). The approximation is wrong: $(1-K^{-1})^I\approx e^{-I/K}$; for $I=1000$, $K=100$ the exact value is $4.3\times10^{-5}$ and $e^{-I/K}=4.5\times10^{-5}$, against $e^{-K/I}=0.905$. Interpretation, checked by simulation: treating the $K$ non-overlapping outcomes as the effective units, almost all of them appear in a bootstrap sample, but its $I$ rows hold only about $K(1-e^{-I/K})$ distinct units, $0.100$ per row against $0.632$ under IID draws. The conclusion that assuming IID draws oversamples survives; the formula does not. The legacy notes also wrote ``overlapping'' for non-overlapping, used $\bar u^{(2)}_j$ in its own normaliser in place of $\bar u^{(2)}_k$, and, like the source, omitted the superscript in $u^{(2)}_{t,j}$ inside $\bar u^{(2)}_j$.

\qffield{See also} \hyperref[def:label-concurrency-uniqueness]{label concurrency and uniqueness}; \hyperref[meth:sample-weights-attribution-decay]{sample weights}; \hyperref[pit:parameter-overfitting]{parameter overfitting}.

\qffield{Further reading} Cross-validation under overlapping labels, with purging and embargo (López de Prado, chap.~7), is a backlog item of the pathway; Rao, Pathak and Koltchinskii (1997) is the sequential resampling scheme López de Prado distinguishes from his (p.~63), not read. The legacy file is kept in \nolinkurl{notes/artefacts/statistical-learning-foundations/}.

\qffield{Index} bootstrap; bagging; out-of-bag accuracy; redundancy; oversampling; average uniqueness; sequential bootstrap; draw probabilities; max\_samples; Monte Carlo

\qffield{Status} Verification: partly checked 2026-10-16 (against the text layer and page images of a local copy of \texttt{alpha-afml}, chap.~4, pp.~62--68; toy example, Monte Carlo and the overlap approximation recomputed). Stability: draft. Currency: timeless.
\end{qfmethod}

\begin{qfmethod}[Sample weights for overlapping labels: return attribution, time decay on cumulative uniqueness, and class weights]\label{meth:sample-weights-attribution-decay}
\qffield{Statement} \emph{Return attribution.} For labels that depend on the sign of the return ($\pm1$, or $\{0,1\}$ for meta-labels), López de Prado weights observation $i$ by the absolute sum of the returns attributed to it,
\[ \tilde w_i=\Bigl|\sum_{t=t_{i,0}+1}^{t_{i,1}}\frac{r_{t-1,t}}{c_t}\Bigr|\ \text{with}\ r_{t-1,t}=\ln(p_t/p_{t-1}),\qquad w_i=\tilde w_i\,\frac{I}{\sum_{j=1}^I\tilde w_j},\qquad \sum_iw_i=I, \]
(log returns, ``so that they are additive'', in Snippet~4.10, unlike the simple return used to define concurrency), because overlapping outcomes would otherwise be over-weighted and ``labels associated with large absolute returns should be given more importance than labels with negligible absolute returns''; the weights sum to $I$ because libraries assume a default weight of 1 \citep[chap.~4, sec.~4.6, pp.~68--69, Snippet~4.10]{alpha-afml}. \emph{Time decay.} Weights are multiplied by $d[x]=\max\{0,a+bx\}$ on $x\in[0,\sum_i\bar u_i]$, cumulative uniqueness rather than calendar time, with $d[\sum_i\bar u_i]=1$ for the newest observation; for $c\in[0,1]$, $d[0]=c$, giving $b=(1-c)/\sum_i\bar u_i$; for $c\in(-1,0)$, $d[-c\sum_i\bar u_i]=0$, giving $b=[(c+1)\sum_i\bar u_i]^{-1}$; in both cases $a=1-b\sum_i\bar u_i$ (sec.~4.7, p.~70). \emph{Class weights} correct for under-represented labels, ``particularly critical'' when the important classes are rare, such as liquidity crises; López de Prado's default is \texttt{class\_weight='balanced'}, or \texttt{'balanced\_subsample'} for bagging (sec.~4.8, pp.~71--72).

\qffield{Conditions} Return attribution ``will not work if there is a `neutral' (return below threshold) case'', for which lower returns would need higher weights; López de Prado advises dropping neutral cases (p.~69). Decay runs on cumulative uniqueness because chronological decay ``would reduce weights too fast in the presence of redundant observations'' (p.~70). Cases: $c=1$ no decay; $0<c<1$ linear decay with every weight positive; $c=0$ decay to zero; $c<0$ zero weight for the oldest observations, those with $x\le-c\sum_i\bar u_i$, a share $|c|$ of cumulative uniqueness (pp.~70--71). In scikit-learn, class weights that do not sum to the number of classes act like a change in the regularisation parameter (p.~71).

\qffield{Practice} Checked (\nolinkurl{afml-checks.py}): attributed weights normalised by $I/\sum_j\tilde w_j$ sum to $I$; the decay constants meet $d[\sum\bar u]=1$ and $d[0]=c$ for $c\ge0$, and $d[-c\sum\bar u]=0$ for $c<0$, for $c=1,0.5,0,-0.5$. Interpretation: López de Prado's implementation evaluates $d$ at the cumulative sums $\bar u_1,\bar u_1+\bar u_2,\dots$ (Snippet~4.11), so the oldest observation sits at $x=\bar u_1>0$ and receives slightly more than $c$; and his printed sum runs from $t=t_{i,0}$ (p.~69), as does Snippet~4.10, which adds the return ending at the event's own bar; the lower limit $t_{i,0}+1$ above sums exactly the returns the label uses, as in \hyperref[def:label-concurrency-uniqueness]{label concurrency}. On a toy path the two give attribution sums $(0.0247,0.0145,0.0190)$ and $(0.0149,0.0146,0.0288)$ for López de Prado's three labels (\nolinkurl{slice32-followup-checks.py}).

\qffield{Pitfalls} The legacy notes dropped the factor $I$ (their weights summed to 1 while claiming $I$), and their ``lower returns should be assigned higher weights'' reversed the rule by dropping its neutral-class context. Two slips are López de Prado's: the text says $d[1]=c$ where its own boundary condition, $a+b\cdot0=c$, gives $d[0]=c$; and ``the oldest portion $cT$ of the observations receive zero weight'' should be a share $-c$ of cumulative uniqueness, not of $T$ (pp.~70--71). Weights built from $t_{i,1}$, $c_t$ and realised returns are functions of future outcomes and belong to training only; see \hyperref[pit:label-weight-leakage]{time-validity of labels and weights}.

\qffield{See also} \hyperref[def:label-concurrency-uniqueness]{label concurrency and uniqueness}; \hyperref[meth:sequential-bootstrap]{sequential bootstrap}; \hyperref[meth:meta-labelling]{meta-labelling}.

\qffield{Further reading} King and Zeng (2001) on rare events, cited by López de Prado (p.~71), not read. The legacy file is kept in \nolinkurl{notes/artefacts/statistical-learning-foundations/}.

\qffield{Index} sample weights; return attribution; absolute return weighting; neutral labels; time decay; cumulative uniqueness; piecewise-linear decay; class weights; balanced class weights; rare events

\qffield{Status} Verification: partly checked 2026-10-16 (against the text layer and page images of a local copy of \texttt{alpha-afml}, chap.~4, pp.~68--72, Snippets~4.10--4.11; weights and decay constants recomputed; the summation's lower limit corrected 2026-10-16). Stability: draft. Currency: timeless; scikit-learn behaviour as described in 2018.
\end{qfmethod}

\begin{qfpitfall}[Time-validity of labels, sample weights and fitted transforms: what is known when, and what may be used only in training]\label{pit:label-weight-leakage}
\qffield{Statement} Several of López de Prado's constructions use information from after the observation they describe. A triple-barrier label and its return use prices up to $t_{i,1}$ \citep[chap.~3, p.~45]{alpha-afml}; average uniqueness ``requires information that is not available until a future time'' and is safe only because it is ``used on the training set \dots\ and not on the testing set'' (chap.~4, p.~62); attribution weights use realised returns over the label's life (p.~69). These are legitimate as training targets and weights. Interpretation, derived here: they leak when (i) a training label's span reaches into the test period, so the training set contains test-period prices; (ii) concurrency $c_t$ is computed with labels from both sets; (iii) a quantity fitted on the full sample (a CUSUM threshold, a barrier multiple, a fractional-differencing order $d$, an EWMA warm-up, a meta-label primary model) is then used in a backtest of that same sample; or (iv) weights or labels are used as features.

\qffield{Conditions} The decision-time inputs, by contrast, are causal as López de Prado defines them: bar thresholds use EWMAs of prior bars (chap.~2, pp.~29--32), and the barrier width uses volatility up to the event (chap.~3, pp.~44--45). His fractional-differencing order is chosen on the whole E-mini history (chap.~5, Snippet~5.4, p.~85), and his quantamental recipe refits the secondary model ``on the same training set'' as the primary (p.~53); both need re-estimation inside each training window when the result is to be evaluated out of sample.

\qffield{Practice} Interpretation, a checklist for one train/test split: compute labels, spans, concurrency, uniqueness and weights from training data; drop (purge) training observations whose spans overlap the test period, and leave a gap (embargo) after it; fit $d$, thresholds and barrier multiples on training data and apply them unchanged to the test period; build meta-labels from out-of-sample primary predictions; and record each choice with its data window.

\qffield{Pitfalls} The legacy notes reported ``no information leakage'' for uniqueness without its training-set condition, and kept the fractional-differencing order as a property of the series rather than an estimate. Purging and embargo are López de Prado's chap.~7, not read when this entry was written; the description above is the entry's. Chap.~7 was read on 2026-10-17 for \hyperref[meth:purged-embargoed-cross-validation]{purged and embargoed cross-validation}, which states his rules; they agree with the checklist's purge and embargo, applied to every split.

\qffield{See also} \hyperref[def:label-concurrency-uniqueness]{label concurrency and uniqueness}; \hyperref[meth:sample-weights-attribution-decay]{sample weights}; \hyperref[meth:meta-labelling]{meta-labelling}; \hyperref[meth:fixed-width-fractional-differencing]{fixed-width fractional differencing}; \hyperref[meth:strategy-production-lifecycle]{strategy production lifecycle} (embargo before live trading, a different use of the word); \hyperref[pit:parameter-overfitting]{parameter overfitting}.

\qffield{Further reading} The pathway's readings on information timing, beginning with \texttt{alpha-jansen} and the point-in-time language-model papers, and López de Prado's chap.~7 are in the backlogs of \texttt{financial-ml-evaluation} and \texttt{statistical-learning-foundations}.

\qffield{Index} look-ahead bias; leakage; label span; training set; test set; purging; embargo; out-of-sample predictions; parameter estimation window; time-valid inputs

\qffield{Status} Verification: partly checked 2026-10-16 (the López de Prado passages cited, chaps.~2--5, as for the linked entries; the checklist is the entry's own and unchecked against chap.~7). Stability: draft. Currency: timeless.
\end{qfpitfall}

\begin{qfmethod}[Purged and embargoed $k$-fold cross-validation: dropping training labels that overlap the test set, a gap after it, and why shuffled $k$-fold and out-of-bag scores are inflated]\label{meth:purged-embargoed-cross-validation}
\qffield{Statement} López de Prado gives two reasons why $k$-fold cross-validation ``fails in finance'': observations ``cannot be assumed to be drawn from an IID process'', and the test set is reused while a model is developed, ``leading to multiple testing and selection bias'' \citep[chap.~7, sec.~7.3, p.~104]{alpha-afml}. Chapter~7 treats only the first. With a serially correlated feature, $X_t\approx X_{t+1}$, and labels formed on overlapping data, $Y_t\approx Y_{t+1}$, putting $t$ and $t+1$ in different sets leaks information, so a classifier can score well ``even if $X$ is an irrelevant feature''; leakage requires $(X_i,Y_i)\approx(X_j,Y_j)$, and ``it does not suffice that $X_i\approx X_j$ or even $Y_i\approx Y_j$'' (p.~105). The remedy has two parts (sec.~7.4, pp.~105--108):
\begin{itemize}
\item \emph{Purging}: with each label a function of prices on a closed interval, $Y_j=f[[t_{j,0},t_{j,1}]]$, drop from the training set every observation $i$ whose label overlaps a test label, that is, if $t_{j,0}\le t_{i,0}\le t_{j,1}$, or $t_{j,0}\le t_{i,1}\le t_{j,1}$, or $t_{i,0}\le t_{j,0}\le t_{j,1}\le t_{i,1}$ (p.~106, Snippet~7.1);
\item \emph{Embargo}: also drop training observations that start within $h$ after the test set, $t_{j,1}\le t_{i,0}\le t_{j,1}+h$, implemented by extending each test label to $t_{j,1}+h$ before purging; training labels that end before the test set begins need no embargo, since they ``contain information that was available at the testing time''; ``a small value $h\approx.01T$ often suffices'' (p.~107, Snippet~7.2).
\end{itemize}
Purging and embargo are to be applied ``whenever we produce a train/test split, whether it is for hyper-parameter fitting, backtesting, or performance evaluation''; Snippet~7.3's \texttt{PurgedKFold} assumes contiguous, unshuffled test blocks (sec.~7.4.3, p.~109).

\qffield{Notation} $T$ is the number of bars; $[t_{i,0},t_{i,1}]$ the span of the prices that determine label $i$ (for a triple-barrier label, from the event to the first barrier touched; \hyperref[meth:triple-barrier-labelling]{triple-barrier labelling}); $h$ the embargo length in bars, \texttt{pctEmbargo}$\,=h/T$ in the snippets.

\qffield{Conditions} The three overlap conditions are equivalent to the two closed intervals intersecting (own check, \nolinkurl{notes/artefacts/statistical-learning-foundations/batch40-cv-checks.py}). The text states the purge rule as dropping $i$ when the information sets satisfy ``$\Phi_i\cap\Phi_j=\emptyset$'' (p.~106, read as a page image at 300~dpi); the next sentence and the snippet drop overlapping labels, so the condition meant is $\Phi_i\cap\Phi_j\neq\emptyset$, an apparent misprint in the 2018 printing read. In the class (p.~109) the left training set keeps only observations whose label ends by the start of the test block, and the right set starts $h$ bars after the first observation that begins after the last test label ends; the same check reproduces these index sets on a toy series, and shows that the class keeps a training label that ends exactly at the test block's first bar, or starts exactly at the last test label's end, which the closed-interval rule would purge (one shared bar); and that its embargo counts observations, \texttt{int(X.shape[0]*pctEmbargo)}, which equals bars only when there is one observation per bar. A test of whether leakage remains: with leakage, performance rises merely as $k\to T$; if, beyond some $k^*$, it stops improving, ``the backtest is not profiting from leaks'' (pp.~106--107).

\qffield{Practice} \emph{Out-of-bag scores.} For bagged trees and random forests on overlapping labels, ``out-of-bag accuracy will be inflated'' because sampling with replacement puts in the training set samples ``very similar to those out-of-bag''; López de Prado advises unshuffled stratified $k$-fold, ignoring out-of-bag accuracy, and a low $k$ (chap.~6, sec.~6.3.3, p.~97), and repeats that random-forest out-of-bag estimates in finance ``are likely to be inflated'' (sec.~6.4, p.~98). \emph{When ordinary $k$-fold is valid.} Bergmeir, Hyndman and Koo show that for a purely autoregressive model, $y_t=g(x_t;\theta)+e_t$ with $x_t$ the lags, standard $K$-fold over the rows of the embedded matrix is valid ``as long as the models considered have uncorrelated errors'', as when the model nests an appropriate one; underfitting should be ruled out first by testing residuals for serial correlation, and in their simulations $K$-fold beat out-of-sample evaluation and non-dependent ($h$-block) cross-validation \citep[working paper of 23 July 2017, abstract and secs.~1--2, pp.~1--3; sec.~7, p.~14]{rf-bergmeir-hyndman-koo-2018}. Cerqueira, Torgo and Mozetič find cross-validation applicable to stationary series, while for real series with non-stationary variation ``the most accurate estimates are produced by out-of-sample methods that preserve the temporal order of observations'' \citep[arXiv v1, 28 May 2019, abstract, p.~1]{rf-cerqueira-torgo-mozetic-2020}. Interpretation: the conditions differ. Bergmeir et al.'s result concerns one-step autoregressive rows with uncorrelated errors; López de Prado's labels span many bars, so neighbouring rows share label prices and the errors of a model that ignores the overlap are not uncorrelated, which is the case their result excludes. Purging removes that overlap but does not address non-stationarity, Cerqueira et al.'s reason for preferring time-ordered holdouts.

\qffield{Pitfalls} Purging and embargo address leakage, not the second failure: choosing among many configurations on the same folds still needs a correction for the number of trials (chap.~7, p.~104; \hyperref[meth:deflated-sharpe-ratio]{deflated Sharpe ratio}). Purged $k$-fold trains on data after the test block, so it is not a walk-forward backtest; López de Prado treats backtesting separately (chaps.~10--16, not read). Cerqueira et al.\ were read in the arXiv preprint (62 real series and three synthetic ones), not the 2020 \emph{Machine Learning} version, which the catalogue describes with 174 real series; Bergmeir et al.\ in the accepted working paper, not the journal version.

\qffield{See also} \hyperref[pit:label-weight-leakage]{time-validity of labels and weights}; \hyperref[def:label-concurrency-uniqueness]{label concurrency and uniqueness}; \hyperref[meth:sequential-bootstrap]{sequential bootstrap} (the bagging remedy, chap.~4); \hyperref[pit:parameter-overfitting]{parameter overfitting}.

\qffield{Further reading} The rest of chap.~7 (sec.~7.5 on scikit-learn's scoring bugs, pp.~109--110, and the exercises) and chap.~12's combinatorial purged cross-validation were not read.

\qffield{Index} cross-validation; $k$-fold; purged $k$-fold; purging; embargo; leakage; overlapping labels; out-of-bag accuracy; random forest; bagging; walk-forward; blocked cross-validation; $h$-block cross-validation; autoregressive cross-validation; non-stationarity

\qffield{Status} Verification: partly checked 2026-10-17 (against the text layer of a local copy of \texttt{alpha-afml}, Wiley 2018, chap.~7, secs.~7.1--7.4.3, pp.~103--109, with p.~106 read as page images, and chap.~6, secs.~6.3.3--6.4, pp.~97--98; \texttt{rf-bergmeir-hyndman-koo-2018}, Monash working paper dated 23 July 2017, pp.~1--3 and 14; \texttt{rf-cerqueira-torgo-mozetic-2020}, arXiv v1, abstract; the overlap rule and the class's index sets recomputed). Stability: draft. Currency: methods as published 2017--2019.
\end{qfmethod}
